\section{Formalization}\label{sec:formal}

The Lean~4 development that accompanies this paper is in the directory \texttt{lean/} next to it: about $12{,}300$ lines in six libraries over Mathlib, with no dependency on leanVM's own formal development. Appendix~\ref{app:lean} lists, for every statement of the paper with machine-checked content and for every Lean file, what corresponds to what, and quotes the main statements verbatim.

\subsection{What is machine-checked}

\begin{center}
\begin{tabularx}{\textwidth}{@{}lX@{}}
\toprule
library & content\\
\midrule
\texttt{DeepFold} & The evaluation functional (Lemma~\ref{lem:E}), the determinant lemma (Lemma~\ref{lem:L1}), the class trichotomy (Lemma~\ref{lem:L2}), the counting theorem (Theorem~\ref{thm:count}), and the fold round with its claims and the predicate $B$ (Lemma~\ref{lem:F}), with lanes and monomial lane weights.\\
\texttt{Protocol} & Pinning (Lemma~\ref{lem:P}), the potential for one batch and for a schedule (Lemma~\ref{lem:T}), goodness of points and the error of the point message (row M0 of Theorem~\ref{thm:rbr-deep}), the exponent maps $e_j$, and the arithmetic of the parameter claims at level $0$ (Proposition~\ref{prop:params}).\\
\texttt{Rounds} & The transitions M1 to M5 of Theorem~\ref{thm:rbr-deep} in the relativised form Theorem~\ref{thm:K} needs, each as a counting statement and as a density with exactly the stated error; uniqueness of near tuples and the descent step.\\
\texttt{Soundness} & The composition. A definition of round-by-round soundness for a public-coin protocol given as a transition system, with the bound of every prover strategy's winning probability by the sum of the errors; an abstract model of Protocol~\ref{prot:D} with one batch of points per level, over any ladder of levels and with its tail; its state function (the invariant of \texttt{Rounds} together with the potential of \texttt{Protocol}); and the theorem that the model is round-by-round sound with exactly the errors of Theorem~\ref{thm:rbr-deep}, hence sound, and knowledge sound in the relativised form, with their sum. Versions with the hypothesis $\HK$ in place of $\Lb$, with no list hypothesis in the unique-decoding range, and two concrete instances with no hypothesis at all. The existence of prover strategies on every ladder and with every tail, under two conditions on the abstract maps.\\
\texttt{Interleave} & The interleaving reductions: Theorem~\ref{thm:I} with its special cases and closed forms, Proposition~\ref{prop:R2}, the two-lane bound, their forms for Reed-Solomon codes with lanes and values in an extension, and the bridge that feeds a single-lane hypothesis over the small field into Lemma~\ref{lem:F}(1).\\
\texttt{Constructions} & The attack side: subspace polynomials are linearized, the flats lemma (Proposition~\ref{prop:flats}), the pole collapse for the affine and for any polynomial combiner (Proposition~\ref{prop:pole}), and three forms of the families at agreement $5\blen/8$ (Proposition~\ref{prop:five8}).\\
\bottomrule
\end{tabularx}
\end{center}

Every file builds with no \texttt{sorry}, no \texttt{axiom}, no \texttt{admit}, no \texttt{native\_decide}, and with every Lean warning treated as an error. Each library has a file \texttt{Axioms.lean} that prints, under \texttt{\#guard\_msgs}, the axioms of a selection of its statements that includes every theorem cited in this paper (50 in \texttt{DeepFold}, 99 in \texttt{Protocol}, 39 in \texttt{Rounds}, 93 in \texttt{Soundness}, 65 in \texttt{Interleave}, three of them in \texttt{Bridge.lean}, and 79 in \texttt{Constructions}): each depends on Lean's three standard axioms \texttt{propext}, \texttt{Classical.choice}, \texttt{Quot.sound} or a subset, and the build fails if one gains an axiom. The libraries also instantiate main statements at fields of the deployed sizes, in the same files and in \texttt{Constructions/Concrete.lean}: at \texttt{GaloisField 2 128}, and at \texttt{GaloisField 2 64} with a quadratic extension for the interleaving reductions and for the concrete instances of the composed theorem. In each instance every typeclass argument is discharged by Mathlib's instances, the characteristic included, so that ``instantiated at the $128$-bit field'' is literally true; for the round-by-round statement \lean{protocolD_rbr} the hypotheses on the analysis constants are discharged as well, at $\Kz=2^9\q$, $M=2\q$ and $\dsc=2^{-119}$. What an instance leaves as an argument is what its statement is about, such as the ladder with its list hypothesis. These instances are a check that the hypotheses on the field can be met. They fix the fields and not leanVM's parameters: those of the interleaving reductions (\lean{concrete_unique}, \lean{concrete_two_mul}, \lean{concrete_descent} in \leanfile{Interleave/Axioms.lean}) have four lanes and an arbitrary domain, and the instances of the composed theorem in which every hypothesis is discharged are the two concrete ones of \S\ref{sec:composed}, neither of which is a deployed parameter set.

\subsection{The composed theorem, exactly}\label{sec:composed}

The statement a reader can cite is \lean{protocolD_sound} (\leanfile{Soundness/Sound.lean}; verbatim in Appendix~\ref{app:lean}). In words: let the field be finite of characteristic two with $\q$ elements, $\dsc>0$ rational, and $M\ge2\q$, $\Kz$ integers with $\Kz(\Kz-1)(\Kz-M)>0$. Consider the model of Protocol~\ref{prot:D} on any ladder of levels, with any tail and any number $J$ of input claims, each level having at least one fold round, a nonempty domain with $\binom{\blen}{2}\le M$, thresholds that grow by at least one per round, and lists of at most $L$ members on the diagonal of \S\ref{sec:hyp}. Then for every deterministic prover strategy, the probability over the verifier's messages that the verifier accepts while it is not the case that the final polynomial descends to a member of the level-$0$ list that matches the commit-phase answers and satisfies all input claims, is at most the sum of the errors of rows M0 to M5 of Theorem~\ref{thm:rbr-deep} over the messages of the run. It follows from \lean{protocolD_rbr} (\leanfile{Soundness/ProtocolD.lean}), which exhibits a state function with properties (i) to (iv) and those errors, and from the union bound \lean{StateFunction.winProb_le} (\leanfile{Soundness/RBR.lean}). \lean{protocolD_sound_plain} is the same bound for ``the verifier accepts false input claims'', and \lean{protocolD_sound_extract} for ``the verifier accepts and the descent fails''. \lean{protocolD_sound_of_HK} replaces the list hypothesis by $\HK$; \lean{protocolD_sound_unique} has no list hypothesis when every level is in the unique-decoding range; \lean{protocolD_sound_concrete_unique_le} discharges every hypothesis on one concrete level with no tail over a field of $2^{128}$ elements (domain the whole subfield of $2^{64}$ elements, degree bound $2^{24}+5$, six fold rounds, twelve points, $128$ queries, $512$ claims, thresholds $2^{63}+2^{24}+j$ in the unique-decoding range) and bounds every strategy's winning probability by $2^{-109}$. Its purpose is to show that the hypotheses of the theorem can be met together. It is not a deployed parameter set: the rate is about $2^{-40}$, and with one level and no tail the run has no message of rows M3 and M5. \lean{protocolD_sound_concrete_ladder} discharges every hypothesis on a ladder of two concrete levels with a tail of two rounds, so that a run has a message of every row M0 to M5: the first level is the same, the second has the same domain, degree bound $2^{20}+5$, four fold rounds, twelve points, $128$ queries and thresholds $2^{63}+2^{20}+j$, its map $U\mapsto\mpoly_{g_U}$ is the first lane of the tuple, and the tables of the tail are $(\inner{W}{P},0,0,0)$ and $(1,1,1,1)$. Its bound is the sum of the errors of the run, which is not evaluated to a number, and it is not a deployed parameter set either. \lean{err_level0} gives the level-$0$ errors for a list of $2^{40}$ before their grinds (at most $2^{-110}$ for the points, $2^{-119}$ for the batching, $2^{-118}$ for a fold, $2^{-127}$ for a tail round).

The theorem is about a \emph{model}, and the model is abstract in these respects.
\begin{itemize}[itemsep=1pt,topsep=2pt]
\item A committed word is a function sent by the prover, as in the interactive oracle model. There are no Merkle trees.
\item The map $U\mapsto\mpoly_{g_U}$ of each level is any linear map from tuples to polynomials. The high part sent at a level boundary is any polynomial on which the running weight vanishes and which keeps the sum with the image of every list member within the degree bound. The tail is a sumcheck on two tables whose inner product is the running weight applied to the final polynomial. In the protocol these are provided by the novel-basis encoder, by $\Wn{\kappa}H$ with (S1), and by the terminal weight with the low part; that the real encoder has these properties is \emph{not} formalized. Nothing else about the encoder enters the composition.
\item A weighted claim is any linear functional on tuples with a value, and the claims may depend on the level-$0$ points.
\item Errors are those of the interactive protocol with uniform verifier messages, before any grind. Row M1 is $(J-1)/\q$: the ring-switching scalars of leanVM's stacked opening are not modeled.
\item Each level has one batch of points. The ladder schedules of \S\ref{sec:points} are not composed in Lean; for them only the potential lemma and the error of the first point message are checked.
\item The bound is over the prover strategies of the model, and the prover's messages are typed: a high part and a final polynomial must have the properties above. At every other position the prover has a message outright. For abstract data the two typed messages can be missing, and the statement is then true and empty: no final polynomial exists if, say, the tail has no round, both of its tables are the constant one and the running weight is zero (\lean{finalMsg_isEmpty}), and no high part exists if the first threshold of the next level is zero and its map sends a tuple within its degree bound to a polynomial above the degree bound of the level that ends (\lean{handover_isEmpty}). Two conditions exclude this: the tables of the tail at the zero polynomial have inner product zero (\lean{Terminal.AdmitsZero}; it holds with no tail, and whenever the tables represent the running weight), and the map of each later level sends the tuples within its degree bound to polynomials within the degree bound of the level before it (\lean{Level.Fits}, \lean{ladderFits}). Under them prover strategies exist from every position of the protocol, for every ladder, every tail length and every number of rounds (\lean{strategy_nonempty_of_fits}, \lean{strategy_nonempty_ladder} in \leanfile{Soundness/Sound.lean}; verbatim in Appendix~\ref{app:lean}), so on every such ladder the composed theorem is not vacuous; both conditions hold for the two concrete instances (\lean{strategy_nonempty_one_level}, \lean{strategy_nonempty_concrete_ladder}). The strategies exhibited send zeros. That the honest prover of the real encoder is a strategy is completeness and is not formalized.
\end{itemize}

\subsection{How it was built}

The development uses Lean \texttt{v4.36.0} and Mathlib \texttt{v4.36.0}, the versions leanVM's own formal development pins; \texttt{lean/lake-manifest.json} records the revisions. The sources that accompany this paper were built with Lake from an empty build directory, with every Lean warning treated as an error: all six libraries, $18{,}369$ build jobs, Mathlib's included, ending with \texttt{Build completed successfully}. The files that were built are identical to the accompanying ones, file by file, by their SHA-256 digests. A control theorem proved by \texttt{sorry}, added to one of the libraries, is rejected by Lake under the same options. To reproduce, on a machine where fetching the toolchain and Mathlib's cache is acceptable: \texttt{cd lean \&\& lake exe cache get \&\& lake build}. The \texttt{README.md} files of \texttt{lean/} and of the library directories say what each file contains.

\subsection{What the Lean statements are, and are not}

The Lean statements are counting statements over an arbitrary finite field (of characteristic two only where needed): a probability is the cardinality of a set of verifier messages divided by the number of messages, in $\mathbb Q$. Where a Lean statement is weaker than the text or covers part of it, the status line of the statement says so (Lemma~\ref{lem:L2}(a), (b) is checked through its consequence; the one-batch potential of Lemma~\ref{lem:T} uses the indicator of the exceptional set; the lists of \S\ref{sec:limit} are lower bounds). Otherwise the Lean statements differ from the text in ways that make them more general: any $M\ge\max(2\q,\binom\blen2)$; weights as arbitrary linear functionals on tuples; exponents as an arbitrary injective bounded map (with the concrete $e_j$ proved to be one); a finite set of points of any field as the domain.

Not machine-checked:
\begin{itemize}[itemsep=1pt,topsep=2pt]
\item \emph{The hypothesis.} $\Lb$ below the Johnson radius is not a theorem of anyone (\S\ref{sec:lists}); the Lean development takes it, or $\HK$, as a hypothesis, exactly as the paper does, and discharges it only in the unique-decoding range. The Johnson bound and the counting bound are not connected to the composed theorem in Lean, so Corollary~\ref{cor:johnson} is by hand. None of the list bounds of \S\ref{sec:proved} for special classes of words is formalized, and neither are the statements of \S\ref{sec:allwords}.
\item \emph{The link from the model to the protocol}: the identities of Annex~B for the real encoder and weights (the novel basis, $\mpoly_g(x)=\inner{W_x}{g}$ at every $x\in\E$, the low and high parts and (S1), the terminal weight), and that the verifier's folded column value is the folded word.
\item \emph{Completeness.} No honest prover is defined in Lean.
\item \emph{Fiat-Shamir, grinding and Merkle commitments.} That a grind of $b$ bits divides an error by $2^b$, the passage from the interactive statement to the non-interactive one, the binding of the hash commitments and the extraction of the committed words from the adversary's hash queries are inherited from the existing analysis of leanVM, with which the state function is compatible in shape; nobody has instantiated it there.
\item \emph{The outer protocol}: the ring switch, the stacking, whatever happens between the commitment and the opening, the recursion.
\item \emph{The ladder schedules} as a composed protocol, and the composition for a predicate other than the canonical one of \S\ref{sec:state}.
\item \emph{The extractor as an algorithm.} The extraction predicate, the existence and uniqueness of the near tuple at each step, and that it folds back, are checked; that it is found by one-error decoding in polynomial time is not.
\item \emph{Exact counts on the attack side} ($56\blen-392$, $35\blen-245$, the Gaussian binomials). For the families at $5\blen/8$ lower bounds linear in $\blen$ are checked, and for the flats the count of hyperplanes.
\item \emph{Numbers beyond level $0$}: the tables of \S\ref{sec:params} other than the entries of Proposition~\ref{prop:params}.
\item \emph{leanVM's concrete field and domain.} The theorems are quantified over every finite field of characteristic two and every finite additive subgroup, so they apply to $\dom_{\kappa,\rexp}\subseteq\K\subseteq\E$; no isomorphism between Mathlib's \texttt{GaloisField 2 128} and the tower $\K[y]/(y^2+xy+1)$ is stated.
\end{itemize}

\subsection{What formalizing showed}

No error and no missing hypothesis was found in any hand proof that was formalized, and the libraries, written separately, composed with no change to any piece. The process did sharpen several statements. In Lemma~\ref{lem:F} no relation between $\agr$, $\dgr$ and $\blen$ is needed, characteristic two is used in exactly one line (the identity that turns condition (p) into consistency with an answer vector), and the hypothesis can be restricted to pairs with $\deg F<\dgr$. The order of quantifiers is as the protocol needs it: the exceptional set is defined before, and independently of, the answers, the claim, the carried values, the predicate and the prover message, and Lean's binder order makes this a checked fact. Proposition~\ref{prop:R2} holds under $L\le|\K|^s/2$, where the hand statement had $|\K|^s/4$. Theorem~\ref{thm:I} needs $L\ge1$ as written by hand (the Lean recursion truncates). The descent inverts the fold (\S\ref{sec:ks}). Only the transition M3 uses anything of the encoding, and of the novel basis it uses only linearity. The composition added three findings, all reported where they belong in \S\ref{sec:rbr}: read as one protocol, the statement must be the relativised one with the canonical predicate, plain soundness being a corollary; $\Lb$ is used on a diagonal only; and the soundness statement uses neither (S2), nor $\dsc<1$, nor $\agr_i\le\blen_i$. On the attack side, the degree of $Q_V$ in the flats lemma is at most, not exactly, $|V|/2-1$, and a pole that separates the whole list, which the clean count of bad challenges for a pole pair needs, is shown to exist only in a field larger than about $2^{75}$ at $\blen=2^{26}$, so the count applies with poles in $\E$ and not in $\K$.
